%! Boxplots and Comparing Distributions — Unit 1, Lesson 1.6 — Slides (Beamer)
\documentclass[aspectratio=169, 11pt]{beamer}
\usepackage{apstats-beamer}   % provides \navyheader{} and \sectionlabel[color]{}
\usetikzlibrary{positioning}  % -beamer loads tikz but not this library

% The C and Y column types live in apstats-article, which a beamer deck does not
% load — redeclare them here rather than pulling in the article preamble.
\newcolumntype{C}[1]{>{\centering\arraybackslash}p{#1}}
\newcolumntype{Y}{>{\centering\arraybackslash}X}

\begin{document}

% TITLE SLIDE (CourseName is not defined in beamer, so the name is literal)
{
\setbeamercolor{background canvas}{bg=navy}
\begin{frame}[plain]
  \vspace{2.8cm}\hspace{0.55cm}%
  \begin{minipage}{0.9\textwidth}
    {\color{goldacc}\bfseries\small LESSON 1.6}\\[0.5cm]
    {\color{white}\bfseries\LARGE Boxplots and Comparing Distributions}\\[1.2cm]
    {\color{skymid}\small Statistics~$\cdot$~Unit 1}
  \end{minipage}
\end{frame}
}

% GRADUAL RELEASE deck: title → targets → warm-up hook → I DO divider → ONE INSTRUCTION FRAME
% PER NOTES ROW (the five numbers · anatomy · one quarter, the crux · comparing distributions) →
% WE DO (Guided Practice: Two Science Sections) → spoken debrief → close & START THE HOMEWORK,
% which is the individual practice → AP Practice → close.
% There is NO group activity in this course — the notes carry the whole release. There is no
% transfer set and no reflection box; the debrief is spoken and has nothing to fill in.
% Vocabulary is GIVEN up front here, not withheld: the instruction frames ARE the direct
% instruction, and the retired EFFL spoiler rule no longer applies anywhere.
% Homework is due THE FIRST CLASS AFTER TWO STUDY HALLS — never "due next class".
\newcommand{\redink}[1]{{\color{redacc}\bfseries #1}}
% \bxp{y-center}{half-height}{min}{Q1}{median}{Q3}{max} — every display is pre-drawn.
\newcommand{\bxp}[7]{%
  \draw[thick, navy] (#3,#1) -- (#4,#1);
  \draw[thick, navy] (#6,#1) -- (#7,#1);
  \draw[thick, navy] (#3,{#1-#2}) -- (#3,{#1+#2});
  \draw[thick, navy] (#7,{#1-#2}) -- (#7,{#1+#2});
  \draw[thick, navy, fill=skymid] (#4,{#1-#2}) rectangle (#6,{#1+#2});
  \draw[very thick, navy] (#5,{#1-#2}) -- (#5,{#1+#2});}

% ── LEARNING TARGETS ─────────────────────────────────────────────────────────
\begin{frame}[t, plain]
  \navyheader{Today's targets}
  \sectionlabel{Lesson 1.6}
  By the end of today, I can\ldots
  \begin{itemize}\setlength{\itemsep}{5pt}
    \item report the \textbf{five-number summary} and point to each of the five numbers on a
          \textbf{boxplot}.
    \item explain why every section of a boxplot holds the \textbf{same quarter} of the data.
    \item decide with the $1.5 \times$ IQR rule whether a value is an \textbf{outlier}, and say
          how a \textbf{modified boxplot} draws it.
    \item compare two distributions from \textbf{parallel boxplots} --- in context, with
          comparative language.
  \end{itemize}
  \begin{block}{How today works}
    Warm-up $\to$ \textbf{notes together} $\to$ \textbf{one we work as a class} $\to$
    debrief $\to$ \textbf{you start the homework here, alone}.
  \end{block}
\end{frame}

% ── WARM-UP HOOK ─────────────────────────────────────────────────────────────
{
\setbeamercolor{background canvas}{bg=navy}
\begin{frame}[plain]
  \vspace{1.5cm}
  \begin{center}
    {\color{skymid}\bfseries\small ONE PICTURE. FOUR SECTIONS. TWENTY PLAYERS.}\\[0.9cm]
    {\color{white}\bfseries\Large One section is more than five times as wide as
    another.\\[0.4cm] So how many players are in each?}\\[0.9cm]
    {\color{goldacc}\small Do not guess. Count. --- Keep your warm-up numbers where you can
    see them; they are about to get a name.}
  \end{center}
\end{frame}
}

% ── I DO — direct instruction begins ─────────────────────────────────────────
{
\setbeamercolor{background canvas}{bg=navy}
\begin{frame}[plain]
  \vspace{1.8cm}\hspace{0.8cm}%
  \begin{minipage}{0.86\textwidth}
    {\color{goldacc}\bfseries\small GUIDED NOTES \ \textbullet\ \ I DO}\\[0.7cm]
    {\color{white}\bfseries\Large Open to your Guided Notes page. Fill each blank as we
    name it --- not before.}\\[0.9cm]
    {\color{skymid}\small Five terms go in the vocabulary box today. Write each one the
    moment we use it.}
  \end{minipage}
\end{frame}
}

% ── NOTES ROW 1 ──────────────────────────────────────────────────────────────
\begin{frame}[t, plain]
  \navyheader{Five numbers, drawn to scale}
  \sectionlabel{notes row 1, The Five Numbers $\cdot$ problems 1--6}
  \vspace{2pt}
  \begin{columns}[t]
    \begin{column}{0.52\textwidth}
      \begin{center}
      \begin{tikzpicture}[scale=0.9]
        \bxp{0.9}{0.42}{0}{0.9}{1.9}{4.1}{8.0}
        \draw[->, thick] (-0.2,0) -- (8.6,0);
        \foreach \x/\lab in {0/min, 0.9/$Q_1$, 1.9/med, 4.1/$Q_3$, 8.0/max} {
          \draw[linegray, very thin] (\x,0.05) -- (\x,1.32);
          \draw (\x,-0.08) -- (\x,0.08);
          \node[below, font=\tiny] at (\x,-0.1) {\lab};}
      \end{tikzpicture}
      \end{center}
    \end{column}
    \begin{column}{0.44\textwidth}
      \begin{itemize}\setlength{\itemsep}{6pt}
        \item The \textbf{five-number summary}: minimum, $Q_1$, median, $Q_3$, maximum
              (EK UNC-1.L.1). Those are your warm-up numbers.
        \item A \textbf{boxplot} is that summary drawn to scale on a number line ---
              \redink{five numbers, nothing else.}
        \item Twenty free-throw totals: min 30, $Q_1$ 34.5, median 39.5, $Q_3$ 50.5, max 74.
      \end{itemize}
    \end{column}
  \end{columns}
  \vspace{6pt}
  \begin{center}
    {\small\color{slate} Problem~6: which \emph{part} of the plot marks $Q_1$? The box has two
    edges --- say which one.}
  \end{center}
\end{frame}

% ── NOTES ROW 2 ──────────────────────────────────────────────────────────────
\begin{frame}[t, plain]
  \navyheader{Anatomy: box, whiskers, and the rider who was set apart}
  \sectionlabel{notes row 2, Anatomy of a Boxplot $\cdot$ problems 7--12}
  \vspace{2pt}
  \begin{columns}[t]
    \begin{column}{0.52\textwidth}
      \begin{center}
      \begin{tikzpicture}[scale=0.9]
        \bxp{0.9}{0.42}{0.2}{0.9}{1.9}{3.4}{5.0}
        \fill[navy] (7.6,0.9) circle (0.10);
        \draw[->, thick] (-0.2,0) -- (8.6,0);
        \node[font=\tiny, above] at (7.6,1.08) {58 min};
        \node[font=\tiny, above] at (5.0,1.08) {40 min};
      \end{tikzpicture}
      \end{center}
    \end{column}
    \begin{column}{0.44\textwidth}
      \begin{itemize}\setlength{\itemsep}{6pt}
        \item The \textbf{box} runs $Q_1$ to $Q_3$ --- the \textbf{middle 50\%}. The line inside
              is the \textbf{median}.
        \item A \textbf{whisker} reaches the most extreme value that is \redink{not} an outlier
              (EK UNC-1.L.2).
        \item Route B: IQR $= 12$, cut-off $32 + 1.5(12) = \mathbf{50}$, and $58 > 50$. So 58
              gets \textbf{its own dot} and the whisker stops at \textbf{40}.
      \end{itemize}
    \end{column}
  \end{columns}
  \vspace{6pt}
  \begin{center}
    {\small\color{slate} That is a \textbf{modified boxplot} --- the version the AP exam draws.
    Problem~8: how many of the twenty sit \emph{inside} the box?}
  \end{center}
\end{frame}

% ── NOTES ROW 3 — THE CRUX ───────────────────────────────────────────────────
\begin{frame}[t, plain]
  \navyheader{Why the wide section is not the full one}
  \sectionlabel{notes row 3, One Quarter $\cdot$ problems 13--17 --- the whole point of today}
  \vspace{6pt}
  \begin{itemize}\setlength{\itemsep}{8pt}
    \item Twenty players. Four sections. Your problem~13 counts: \textbf{5 \quad 5 \quad 5 \quad
          5}.
    \item That is \textbf{25\% \quad 25\% \quad 25\% \quad 25\%} --- and it is always true, in
          every boxplot ever drawn.
    \item The right section is more than five times as wide as the left. It holds the
          \redink{same five players}, just spread from 52 to 74 instead of packed into 30 to 33.
  \end{itemize}
  \vspace{8pt}
  \begin{block}{Essential Knowledge (UNC-1.L.2)}
    \textbf{Each section of a boxplot holds one quarter of the data.} Width reports how far apart
    those values are --- \textbf{spread}, never \textbf{count}. And problem~17 is the price: a
    gap, a clump, or a second peak all vanish inside the box, so for a shape question ask for a
    \redink{dotplot or histogram} instead.
  \end{block}
\end{frame}

% ── NOTES ROW 4 ──────────────────────────────────────────────────────────────
\begin{frame}[t, plain]
  \navyheader{Comparing two distributions}
  \sectionlabel{notes row 4, Comparing Distributions $\cdot$ problems 18--21}
  \vspace{4pt}
  \begin{columns}[t]
    \begin{column}{0.46\textwidth}
      \begin{block}{The checklist (EK UNC-1.N.1, 1.O.1)}
        On a shared axis they are \textbf{parallel (side-by-side) boxplots}.\\[3pt]
        \textbf{Center} $\cdot$ \textbf{Variability} $\cdot$ \textbf{Unusual features} --- all
        three, every time, each with a \textbf{comparative word} and each \textbf{in context}.
      \end{block}
    \end{column}
    \begin{column}{0.48\textwidth}
      \begin{block}{Mean vs.\ median (EK UNC-1.M.2)}
        Symmetric: the two sit close together.\\[3pt]
        Skewed \textbf{right}: the mean is to the \textbf{right} of the median --- that is
        problem~21, Route B.\\[3pt]
        Skewed \textbf{left}: to the \textbf{left}.
      \end{block}
    \end{column}
  \end{columns}
  \vspace{10pt}
  \begin{center}
    {\large Describing route A, then describing route B, is \redink{not} a comparison ---
    and earns nothing.}
  \end{center}
\end{frame}

% ── WE DO ────────────────────────────────────────────────────────────────────
\begin{frame}[t, plain]
  \navyheader{Guided Practice: Two Science Sections}
  \sectionlabel[goldacc]{We do $\cdot$ problems 22--26 $\cdot$ you hold the pen}
  Two sections of the same course, one test, two boxplots on one axis.
  \begin{itemize}\setlength{\itemsep}{5pt}
    \item Read both medians and both IQRs --- with units, and with the subtraction shown.
    \item Which section is skewed, and which way? Does either show an outlier?
    \item \textbf{Problem 26:} write the comparison as \textbf{one paragraph}, in context.
  \end{itemize}
  \begin{block}{The standard, out loud}
    A paragraph that describes Period 2 and then describes Period 5 earns \redink{nothing}.
    The words that earn the point are \emph{higher than} and \emph{more variable than} --- with
    the variable, its units, and both sections named.
  \end{block}
\end{frame}

% ── DEBRIEF — spoken, nothing to fill in ─────────────────────────────────────
\begin{frame}[t, plain]
  \navyheader{Debrief: all four sections held five}
  \sectionlabel{8 minutes $\cdot$ pens down, papers up --- we say these aloud}
  \vspace{6pt}
  \begin{columns}[t]
    \begin{column}{0.46\textwidth}
      \begin{block}{Nothing was hidden}
        Four sections, 4.5 to 23.5 shots wide, and every one of them held five of the twenty.
      \end{block}
    \end{column}
    \begin{column}{0.46\textwidth}
      \begin{block}{Only the spacing changed}
        The wide section's five are strung from 52 to 74. The narrow one's five are stacked
        inside 30 to 33.
      \end{block}
    \end{column}
  \end{columns}
  \vspace{12pt}
  \begin{block}{Two things to say out loud}
    State the quarter rule \redink{without using the word ``wide''}. Then give me one comparison
    sentence about something that was on no page today --- with a comparative word and a unit in
    it.
  \end{block}
\end{frame}

% ── YOU DO — close & start the homework, the individual practice ──────────────
\begin{frame}[t, plain]
  \navyheader{You do: start the homework}
  \sectionlabel{Last 8 minutes $\cdot$ on your own $\cdot$ finish it at home}
  \begin{itemize}\setlength{\itemsep}{5pt}
    \item Eleven reading times: five numbers, then which of three plots was built from them.
    \item ``The long right whisker means most of the commute times are large.'' Correct that.
    \item One $1.5 \times$ IQR check on a 52-minute bus delay, arithmetic shown.
    \item Two pizza shops: the full comparison, in context.
    \item One multiple-choice item, with a one-sentence justification.
  \end{itemize}
  \begin{block}{Silent and alone --- I am circulating}
    \emph{Item 2} is the one I am reading over your shoulder, and a correction with no
    \redink{number} in it is only half an answer. Do 1, 2, and 3 while I am here; 4 and 5 travel
    home.
  \end{block}
\end{frame}

% ── HOMEWORK (started in class, scored) ─────────────────────────────────────
\begin{frame}[t, plain]
  \navyheader{Homework --- start it now}
  \sectionlabel[goldacc]{5 exercises $\cdot$ scored $\cdot$ due the first class after two study halls}
  \begin{itemize}\setlength{\itemsep}{5pt}
    \item Five numbers from an ordered list --- then which of three plots is it?
    \item The long-right-whisker claim, corrected with a count.
    \item One $1.5 \times$ IQR check, arithmetic shown.
    \item Two pizza shops: the full comparison, in context.
    \item One multiple-choice item, with a one-sentence justification.
  \end{itemize}
  \begin{block}{This one is graded}
    This is your \redink{scored} homework, on paper, due the first class after two study halls.
    Start with 1, 2, and 3 while I am still here to answer questions --- finish the rest at home.
  \end{block}
\end{frame}

% ── AP PRACTICE & PREVIEW (assigned, unscored) ──────────────────────────────
\begin{frame}[t, plain]
  \navyheader{AP Practice \& what's next}
  \sectionlabel[goldacc]{AP Practice --- assigned, not scored}
  \begin{itemize}\setlength{\itemsep}{5pt}
    \item Four multiple-choice items on weekday versus weekend tips.
    \item One free-response set: report, test an outlier, compare, and choose a better display.
    \item Optional: find published side-by-side boxplots and rewrite their comparison sentence.
  \end{itemize}
  \begin{block}{Not graded --- but this is what the exam looks like}
    The AP Practice page is \textbf{not scored}. Work it for the reps, and bring what you get
    stuck on to study hall or office hours. The \redink{graded} homework is the section behind
    it.
  \end{block}
\end{frame}

% ── CLOSE ────────────────────────────────────────────────────────────────────
{
\setbeamercolor{background canvas}{bg=navy}
\begin{frame}[plain]
  \vspace{1.5cm}\hspace{0.8cm}%
  \begin{minipage}{0.86\textwidth}
    {\color{goldacc}\bfseries\small BEFORE YOU GO}\\[0.7cm]
    {\color{white}\bfseries\Large Every section of a boxplot holds the same quarter of the data.
    Width tells you how stretched out those values are --- it never tells you how many there
    are.}\\[0.8cm]
    {\color{skymid}\small Next: Lesson 1.7 --- scores from two completely different tests, and
    a fair way to say which one was better.}
  \end{minipage}
\end{frame}
}

\end{document}
